\documentclass{article}

\def\npart {IA}
\def\nterm {Michaelmas}
\def\nyear {2023}
\def\nlecturer {Prof N.\ Peake}
\def\ncourse {Vector and Matrices}

\usepackage{alltt}
\usepackage{amsmath, amsfonts, amssymb, amsthm}
\usepackage{booktabs}
\usepackage{caption}
\usepackage{enumitem}
\usepackage{fancyhdr}
\usepackage{graphicx}
\usepackage{mathdots}
\usepackage{mathtools}
\usepackage{microtype}
\usepackage{multirow}
\usepackage{pdflscape}
\usepackage{pgfplots}
\usepackage{siunitx}
\usepackage{textcomp}
\usepackage{slashed}
\usepackage{tabularx}
\usepackage{tikz}
\usepackage{tkz-euclide}
\usepackage[normalem]{ulem}
\usepackage[all]{xy}

\pgfplotsset{compat=1.18}
\pagestyle{fancyplain}

\lhead{\emph{\nouppercase{\leftmark}}}
\rhead{
  \ifnum\thepage=1
  \else
    \npart\ \vline\ \ncourse
  \fi}

\def\nauthor{Marcus Ng}
\author{Based on lectures by \nlecturer \\\small Notes taken by \nauthor}
\date{\nterm\ \nyear}
\title{Part \npart\ --- \ncourse}

\newcommand*{\Cdot}{{\raisebox{-0.25ex}{\scalebox{1.5}{$\cdot$}}}}
\setlist[enumerate,1]{label={(\roman*)}}

\newcommand {\pd}[2][ ]{
  \ifx #1 { }
    \frac{\partial}{\partial #2}
  \else
    \frac{\partial^{#1}}{\partial #2^{#1}}
  \fi
}

% Theorems
\theoremstyle{definition}
\newtheorem*{aim}{Aim}
\newtheorem*{axiom}{Axiom}
\newtheorem*{claim}{Claim}
\newtheorem*{cor}{Corollary}
\newtheorem*{conjecture}{Conjecture}
\newtheorem*{defi}{Definition}
\newtheorem*{eg}{Example}
\newtheorem*{ex}{Exercise}
\newtheorem*{fact}{Fact}
\newtheorem*{law}{Law}
\newtheorem*{lemma}{Lemma}
\newtheorem*{notation}{Notation}
\newtheorem*{prop}{Proposition}
\newtheorem*{question}{Question}
\newtheorem*{problem}{Problem}
\newtheorem*{rrule}{Rule}
\newtheorem*{thm}{Theorem}
\newtheorem*{assumption}{Assumption}
\newtheorem*{assert}{Assertion}

\newtheorem*{remark}{Remark}
\newtheorem*{warning}{Warning}
\newtheorem*{exercise}{Exercise}

\newtheorem{nthm}{Theorem}[section]
\newtheorem{nlemma}[nthm]{Lemma}
\newtheorem{nprop}[nthm]{Proposition}
\newtheorem{ncor}[nthm]{Corollary}

\SetEnumitemKey{cases}{
  label=\underline{Case~\arabic*:},
  ref={\arabic*},
  align=left,
  leftmargin=\parindent
}

% \renewcommand{\labelitemi}{--}
% \renewcommand{\labelitemii}{$\circ$}
% \renewcommand{\labelenumi}{(\roman{*})}

\let\stdsection\section
\renewcommand\section{\newpage\stdsection}

% Strike through
\def\st{\bgroup \ULdepth=-.55ex \ULset}


%%%%%%%%%%%%%%%%%%%%%%%%%
%%%%% Maths Symbols %%%%%
%%%%%%%%%%%%%%%%%%%%%%%%%

% Logical symbols
\newcommand{\contradiction}{
  \begin{tikzpicture}[rotate=45,x=0.5ex,y=0.5ex]
  \draw[line width=.1ex] (0,2) -- (3,2) (0,1) -- (3,1) (1,3) -- (1,0) (2,3) -- (2,0);
  \end{tikzpicture}
}


\let\U\relax
\let\C\relax
\let\G\relax

% Matrix groups
% \newcommand{\GL}{\mathrm{GL}}
% \newcommand{\Or}{\mathrm{O}}
% \newcommand{\PGL}{\mathrm{PGL}}
% \newcommand{\PSL}{\mathrm{PSL}}
% \newcommand{\PSO}{\mathrm{PSO}}
% \newcommand{\PSU}{\mathrm{PSU}}
% \newcommand{\SL}{\mathrm{SL}}
% \newcommand{\SO}{\mathrm{SO}}
% \newcommand{\Spin}{\mathrm{Spin}}
% \newcommand{\Sp}{\mathrm{Sp}}
% \newcommand{\SU}{\mathrm{SU}}
\newcommand{\U}{\mathrm{U}}
% \newcommand{\Mat}{\mathrm{Mat}}

% Matrix algebras
% \newcommand{\gl}{\mathfrak{gl}}
% \newcommand{\ort}{\mathfrak{o}}
% \newcommand{\so}{\mathfrak{so}}
% \newcommand{\su}{\mathfrak{su}}
% \newcommand{\uu}{\mathfrak{u}}
% \renewcommand{\sl}{\mathfrak{sl}}

% Special sets
\newcommand{\C}{\mathbb{C}}
\newcommand{\CP}{\mathbb{CP}}
\newcommand{\GG}{\mathbb{G}}
\newcommand{\N}{\mathbb{N}}
\newcommand{\Q}{\mathbb{Q}}
\newcommand{\R}{\mathbb{R}}
\newcommand{\RP}{\mathbb{RP}}
\newcommand{\T}{\mathbb{T}}
\newcommand{\Z}{\mathbb{Z}}
\renewcommand{\H}{\mathbb{H}}

% Brackets
\newcommand{\abs}[1]{\left\lvert #1\right\rvert}
\newcommand{\bket}[1]{\left\lvert #1\right\rangle}
\newcommand{\brak}[1]{\left\langle #1 \right\rvert}
\newcommand{\braket}[2]{\left\langle #1\middle\vert #2 \right\rangle}
\newcommand{\bra}{\langle}
\newcommand{\ket}{\rangle}
\newcommand{\norm}[1]{\left\lVert #1\right\rVert}
\newcommand{\normalorder}[1]{\mathop{:}\nolimits\!#1\!\mathop{:}\nolimits}
\newcommand{\tv}[1]{|#1|}
\renewcommand{\vec}[1]{\boldsymbol{\mathbf{#1}}}

% not-math
% \newcommand{\bolds}[1]{{\bfseries #1}}
% \newcommand{\cat}[1]{\mathsf{#1}}
% \newcommand{\ph}{\,\cdot\,}
% \newcommand{\term}[1]{\emph{#1}\index{#1}}
% \newcommand{\phantomeq}{\hphantom{{}={}}}

% Probability
% \DeclareMathOperator{\Bernoulli}{Bernoulli}
% \DeclareMathOperator{\betaD}{beta}
% \DeclareMathOperator{\bias}{bias}
% \DeclareMathOperator{\binomial}{binomial}
% \DeclareMathOperator{\corr}{corr}
% \DeclareMathOperator{\cov}{cov}
% \DeclareMathOperator{\gammaD}{gamma}
% \DeclareMathOperator{\mse}{mse}
% \DeclareMathOperator{\multinomial}{multinomial}
% \DeclareMathOperator{\Poisson}{Poisson}
% \DeclareMathOperator{\var}{var}
% \newcommand{\E}{\mathbb{E}}
% \newcommand{\Prob}{\mathbb{P}}

% Algebra
% \DeclareMathOperator{\adj}{adj}
% \DeclareMathOperator{\Ann}{Ann}
% \DeclareMathOperator{\Aut}{Aut}
% \DeclareMathOperator{\Char}{char}
% \DeclareMathOperator{\disc}{disc}
% \DeclareMathOperator{\dom}{dom}
% \DeclareMathOperator{\fix}{fix}
% \DeclareMathOperator{\Hom}{Hom}
% \DeclareMathOperator{\id}{id}
% \DeclareMathOperator{\image}{image}
% \DeclareMathOperator{\im}{im}
% \DeclareMathOperator{\re}{re}
% \DeclareMathOperator{\tr}{tr}
% \DeclareMathOperator{\Tr}{Tr}
% \newcommand{\Bilin}{\mathrm{Bilin}}
% \newcommand{\Frob}{\mathrm{Frob}}

% Others
% \newcommand\ad{\mathrm{ad}}
% \newcommand\Art{\mathrm{Art}}
% \newcommand{\B}{\mathcal{B}}
% \newcommand{\cU}{\mathcal{U}}
% \newcommand{\Der}{\mathrm{Der}}
% \newcommand{\D}{\mathrm{D}}
% \newcommand{\dR}{\mathrm{dR}}
% \newcommand{\exterior}{\mathchoice{{\textstyle\bigwedge}}{{\bigwedge}}{{\textstyle\wedge}}{{\scriptstyle\wedge}}}
% \newcommand{\F}{\mathbb{F}}
% \newcommand{\G}{\mathcal{G}}
% \newcommand{\Gr}{\mathrm{Gr}}
% \newcommand{\haut}{\mathrm{ht}}
% \newcommand{\Hol}{\mathrm{Hol}}
% \newcommand{\hol}{\mathfrak{hol}}
% \newcommand{\Id}{\mathrm{Id}}
% \newcommand{\lie}[1]{\mathfrak{#1}}
% \newcommand{\op}{\mathrm{op}}
% \newcommand{\Oc}{\mathcal{O}}
% \newcommand{\pr}{\mathrm{pr}}
% \newcommand{\Ps}{\mathcal{P}}
% \newcommand{\pt}{\mathrm{pt}}
% \newcommand{\qeq}{\mathrel{``{=}"}}
% \newcommand{\Rs}{\mathcal{R}}
% \newcommand{\Vect}{\mathrm{Vect}}
% \newcommand{\wsto}{\stackrel{\mathrm{w}^*}{\to}}
% \newcommand{\wt}{\mathrm{wt}}
% \newcommand{\wto}{\stackrel{\mathrm{w}}{\to}}
% \renewcommand{\d}{\mathrm{d}}
% \renewcommand{\P}{\mathbb{P}}
%\renewcommand{\F}{\mathcal{F}}

\let\Im\relax
\let\Re\relax

% \DeclareMathOperator{\area}{area}
% \DeclareMathOperator{\card}{card}
% \DeclareMathOperator{\ccl}{ccl}
% \DeclareMathOperator{\ch}{ch}
% \DeclareMathOperator{\cl}{cl}
% \DeclareMathOperator{\cls}{\overline{\mathrm{span}}}
% \DeclareMathOperator{\coker}{coker}
% \DeclareMathOperator{\conv}{conv}
% \DeclareMathOperator{\cosec}{cosec}
% \DeclareMathOperator{\cosech}{cosech}
% \DeclareMathOperator{\covol}{covol}
% \DeclareMathOperator{\diag}{diag}
% \DeclareMathOperator{\diam}{diam}
% \DeclareMathOperator{\Diff}{Diff}
% \DeclareMathOperator{\End}{End}
% \DeclareMathOperator{\energy}{energy}
% \DeclareMathOperator{\erfc}{erfc}
% \DeclareMathOperator{\erf}{erf}
% \DeclareMathOperator*{\esssup}{ess\,sup}
% \DeclareMathOperator{\ev}{ev}
% \DeclareMathOperator{\Ext}{Ext}
% \DeclareMathOperator{\fst}{fst}
% \DeclareMathOperator{\Fit}{Fit}
% \DeclareMathOperator{\Frac}{Frac}
% \DeclareMathOperator{\Gal}{Gal}
% \DeclareMathOperator{\gr}{gr}
% \DeclareMathOperator{\hcf}{hcf}
\DeclareMathOperator{\Im}{Im}
% \DeclareMathOperator{\Ind}{Ind}
% \DeclareMathOperator{\Int}{Int}
% \DeclareMathOperator{\Isom}{Isom}
% \DeclareMathOperator{\lcm}{lcm}
% \DeclareMathOperator{\length}{length}
% \DeclareMathOperator{\Lie}{Lie}
% \DeclareMathOperator{\like}{like}
% \DeclareMathOperator{\Lk}{Lk}
% \DeclareMathOperator{\Maps}{Maps}
% \DeclareMathOperator{\orb}{orb}
% \DeclareMathOperator{\ord}{ord}
% \DeclareMathOperator{\otp}{otp}
% \DeclareMathOperator{\poly}{poly}
% \DeclareMathOperator{\rank}{rank}
% \DeclareMathOperator{\rel}{rel}
% \DeclareMathOperator{\Rad}{Rad}
\DeclareMathOperator{\Re}{Re}
% \DeclareMathOperator*{\res}{res}
% \DeclareMathOperator{\Res}{Res}
% \DeclareMathOperator{\Ric}{Ric}
% \DeclareMathOperator{\rk}{rk}
% \DeclareMathOperator{\Rees}{Rees}
% \DeclareMathOperator{\Root}{Root}
% \DeclareMathOperator{\sech}{sech}
% \DeclareMathOperator{\sgn}{sgn}
% \DeclareMathOperator{\snd}{snd}
% \DeclareMathOperator{\Spec}{Spec}
% \DeclareMathOperator{\spn}{span}
% \DeclareMathOperator{\stab}{stab}
% \DeclareMathOperator{\St}{St}
% \DeclareMathOperator{\supp}{supp}
% \DeclareMathOperator{\Syl}{Syl}
% \DeclareMathOperator{\Sym}{Sym}
% \DeclareMathOperator{\vol}{vol}
\numberwithin{equation}{section}

\begin{document}
\maketitle{
    \small 
    \noindent\textbf{Complex numbers}\\
    Review of complex numbers, including complex conjugate, inverse, modulus, argument and Argand diagram. Informal treatment of complex logarithm, $n$-th roots and complex powers. de Moivre's theorem.\hspace*{\fill}[2]

    \vspace{10pt}
    \noindent\textbf{Vectors}\\
    Review of elementary algebra of vectors in $\R^3$, including scalar product. Brief discussion of vectors in $\R^n$ and $\C^n$; scalar product and the Cauchy-Schwarz inequality. Concepts of linear span, linear independence, subspaces, basis and dimension.

    \vspace{5pt}
    \noindent Suffix notation: including summation convention, $\delta_{ij}$ and $\varepsilon_{ijk}$. Vector product and triple product: definition and geometrical interpretation. Solution of linear vector equations. Applications of vectors to geometry, including equations of lines, planes and spheres.\hspace*{\fill}[5]

    \vspace{10pt}
    \noindent\textbf{Matrices}\\
    Elementary algebra of $3\times 3$ matrices, including determinants. Extension to $n\times n$ complex matrices. Trace, determinant, non-singular matrices and inverses. Matrices as linear transformations; examples of geometrical actions including rotations, reflections, dilations, shears; kernel and image.\hspace*{\fill}[4]

    \vspace{5pt}
    \noindent Simultaneous linear equations: matrix formulation; existence and uniqueness of solutions, geometric interpretation; Gaussian elimination.\hspace*{\fill}[3]

    \vspace{5pt}
    \noindent Symmetric, anti-symmetric, orthogonal, hermitian and unitary matrices. Decomposition of a general matrix into isotropic, symmetric trace-free and antisymmetric parts.\hspace*{\fill}[1]

    \vspace{10pt}
    \noindent\textbf{Eigenvalues and Eigenvectors}\\
    Eigenvalues and eigenvectors; geometric significance.\hspace*{\fill}[2]

    \vspace{5pt}
    \noindent Proof that eigenvalues of hermitian matrix are real, and that distinct eigenvalues give an orthogonal basis of eigenvectors. The effect of a general change of basis (similarity transformations). Diagonalization of general matrices: sufficient conditions; examples of matrices that cannot be diagonalized. Canonical forms for $2 \times 2$ matrices.\hspace*{\fill}[5]

    \vspace{5pt}
    \noindent Discussion of quadratic forms, including change of basis. Classification of conics, cartesian and polar forms.\hspace*{\fill}[1]

    \vspace{5pt}
    \noindent Rotation matrices and Lorentz transformations as transformation groups.\hspace*{\fill}[1]}

\tableofcontents

%%%%%%%%%%%%%%%%%%%%%%
% Lecture 1: 6/10/2023
%%%%%%%%%%%%%%%%%%%%%%
\section{Complex Numbers}
\subsection{Basic properties}
\begin{defi}[Complex number]
    A \emph{complex number} $z \in \C$ is
    \[
        z = a + ib, \quad a, b \in \R, \quad i = \sqrt{-1}  
    \]
    We define two other functions, $\Re()$ and $\Im()$ such that,
    \[
        \Re(z) = a \quad \Im(z) = b 
    \]
    Clearly, we may note that $\R \subset \C$ in the case where $b = 0$
\end{defi}

We can define algebraic manipulation over the complex numbers,
\begin{align*}
    z_1 \pm z_2 &= a_1 \pm a_2 + i(b_1 \pm b_2) \\
    z_1z_2 &= a_1a_2 - b_1b_2 + i(a_1b_2 + a_2b_1)
\end{align*}
We can define the inverse of a complex number
\begin{equation}
    z^{-1} = \frac{1}{a + ib} = \frac{a - ib}{a^2 + b^2}\label{eq:1}
\end{equation}
this demonstrates that $\C$ is \emph{closed} under addition and multiplication. We also define a complex conjugate of $z = a + ib$ as $\overline{z}$ or $z^*$.
\[
    \overline{z}, z^* = a - ib
\]
The modulus of $z$ as $\abs{z}$
\begin{align*}
    \abs{z} &= \abs{a + ib} = \sqrt{a^2 + b^2} \\
    z\overline{z} &= a^2 + b^2 = \abs{z}^2 \\
    z^{-1} &= \frac{\overline{z}}{\abs{z}^2} \tag{see \eqref{eq:1}}
\end{align*}
We also introduce the Argand Diagram, where the coordinate axes are the Real and Imaginary axes. We can see that $\overline{z}$ is the reflection of $z$ in the $x$-axis.

% draw an argand diagram in tikz
\begin{defi}[Triangle inequality]
    For $z_1, z_2 \in \C$, the triangle inequality states that
    \begin{equation}
        \abs{z_1 + z_2} \leq \abs{z_1} + \abs{z_2}\label{eq:2}
    \end{equation}
\end{defi}

\begin{prop}
    The triangle inequality has another common alternative form,
    \begin{equation}
        \abs{z_1 - z_2} \geq \abs{\abs{z_1} - \abs{z_2}}\label{eq:3}
    \end{equation}
\end{prop}

\begin{proof}
    Write $z_1 = z^{\prime}_1 - z^{\prime}_2$, and $z_2 = z^{\prime}_2$ and substitute them into \eqref{eq:2}. Then,
    \begin{align*}
        &\abs{z^{\prime}_1} \leq \abs{z^{\prime}_1 - z^{\prime}_2} + \abs{z^{\prime}_2} \\
        \Rightarrow& \abs{z^{\prime}_1} - \abs{z^{\prime}_2} \leq\abs{z^{\prime}_1 - z^{\prime}_2}\\
        \Rightarrow& \abs{z^{\prime}_1 - z^{\prime}_2} \geq \abs{z^{\prime}_1} - \abs{z^{\prime}_2} \tag{$\star$}
    \end{align*}
    Notice that our labelling of $z_1$ and $z_2$ was arbitrary, and so we can freely swap the assignment. This gives,
    \[
        \abs{z^{\prime}_1 - z^{\prime}_2} \geq \abs{z^{\prime}_2} - \abs{z^{\prime}_1} \tag{$\star\star$}
    \]
    Since $\abs{z^{\prime}_1 - z^{\prime}_2} \geq (\star) \ \text{and} \ (\star\star)$, hence, \eqref{eq:3}
\end{proof}

\begin{defi}[Polar Representation]
    For a complex number $z = a + ib$ we use the argand diagram to write $z = r(\cos{\theta} + i\sin{\theta})$, where
    \[
        r = \abs{z} \quad \text{and} \quad \theta = \arctan\left(\frac{b}{a}\right) = arg(z)  
    \]
\end{defi}

\begin{remark}
    Note that the mapping $(r, \theta) \rightarrow z$ is unique. However, the reverse $r \rightarrow (r, \theta)$ is \emph{NOT} unique since for each value of $\theta$, $\theta + 2n\pi$ is equally valid.

    In order to make this mapping single valued, we will by convention restrict $\theta$ to the principle value, where
    \[
        -\pi < \theta \leq \pi  
    \]

    This multivalued property becomes quite interesting when we consider complex functions. We can imagine taking the multivalued imaginary $\ln$ function as a multistory carpark, walking round and round up and down. To fix this what we do is to build a brick wall on the upramp (``branch cut'').
\end{remark}

Using complex numbers in this form allows us to view complex arithmetic geometrically. For instance, let
\begin{align*}
    z_k &= x_k + iy_k \quad \text{for} \quad k = 1, 2 \\
    &= r_k(\cos{\theta_k} + i\sin(\theta_k))
\end{align*}
\begin{align*}
    z_1z_2 &= r_1(\cos{\theta_1} + i\sin{\theta_1})r_2(\cos{\theta_2} + i\sin{\theta_2}) \\
    &= r_1r_2[(\cos{\theta_1}\cos{\theta_2} - \sin{\theta_1}\sin{\theta_2}) + i(\cos{\theta_1}\sin{\theta_2} + \sin{\theta_1}\cos{\theta_2})] \\
    &= r_1r_2(\cos{(\theta_1 + \theta_2)} + i\sin{(\theta_1 + \theta_2)}) \numberthis \label{eq:4}
\end{align*}

\subsection{Complex exponential function}
\begin{defi}[Complex exponential function]
    \begin{align*}
        exp(z) &= 1 + z + \frac{z^2}{2!} + \cdots + \frac{z^n}{n!} + \cdots \\
        &=\sum_{n=0}^{\infty}{\frac{z^n}{n!}} \numberthis \label{eq:5}
    \end{align*}
\end{defi}
This function is ``uniformly convergent'' (something that we will cover next term in analysis). That means that we are able to do lots of algebraic manipulation without too much hesitation. =

\begin{prop}[Multiplication]
    \[
        exp(z_1)exp(z_2) = exp(z_1 + z_2)  
    \]
\end{prop}
\begin{proof}
    We will consider the double sum
    \begin{align*}
        \sum_{m,n=0}^{\infty}{a_mn} &= a_{00} + a_{01} + a_{02} + a_{03} + \cdots \\
         &+ a_{10} + a_{11} + a_{12} + \cdots \\
         &+ a_{20} + a_{21} + a_{22} + \cdots \\
         &+ a_{30} + a_{31} + a_{32} + \cdots \\
         &\vdots
    \end{align*}
    Then by collecting terms along the diagonal we can regroup the sum into
    \[
        =\sum_{r=0}^{\infty}\sum_{m=0}^{r}{a_{r-m, m}}
    \]
    Hence, 
    \begin{align*}
        exp(z_1) \cdot exp(z_2) &= \sum_{r=0}^{\infty}\sum_{m=0}^{r}{\frac{z_1^{r-m}}{(r-m)!}\frac{z_2^{m}}{m!}} \\
        &= \sum_{r=0}^{\infty}{\frac{1}{r!}}\sum_{m=0}^{r}{\frac{r!}{(r-m)!m!}z_1^{r-m}z_2^{m}} \\
        &= \sum_{r=0}^{\infty}{\frac{1}{r!} (z_1 + z_2)^r} \\
        &= exp(z_1 + z_2)
    \end{align*}
\end{proof}

%%%%%%%%%%%%%%%%%%%%%%
% Lecture 2: 8/10/2023
%%%%%%%%%%%%%%%%%%%%%%
\end{document}